\section{Reproducibility and Release Provenance}\label{sec:reproducibility}

Reproduction is defined by a source revision, locked dependency graph, policy
document, classifier registry, corpus manifest, external-tool set, execution
environment, and machine-readable result directory. The Rust toolchain is
pinned to version 1.96.1 with a minimal profile plus formatting and linting
components. \texttt{Cargo.lock} fixes Rust dependency resolution.
\texttt{Cargo.toml} fixes unwinding panics in development and release because
panic containment is part of the runtime contract. Rebuilding with another
toolchain, unlocked dependencies, altered panic behavior, or an unrecorded
policy is a different experimental condition.

The development repository is publicly available.
{\raggedright
\url{https://github.com/oleksandrmelnychenko/ecliptix-file-defender}\par}
The clean implementation snapshot inspected for this manuscript revision is
\texttt{0bffc39314bc8ee8\allowbreak 40aa6a85e69f108e\allowbreak 667fa5fd}.
This identifier records the
local source state but is not substituted for the commit field of a future
controlled run. Submission still requires an immutable archive, its digest,
and a persistent archive identifier. Until those exist, the repository is an
inspectable development artifact rather than an exactly reproducible published
artifact.

The ordinary quality workflow checks formatting, builds and tests every
workspace target and feature, rejects lint warnings, packages every crate, and
runs dependency advisory and license/source policy gates. Research scripts are
syntax checked and JSON inputs are parsed. A macOS job rebuilds the Apple
XCFramework, compares the generated distribution, header, and Swift package
metadata against the committed artifact, then builds and tests the Swift
package. A benchmark smoke job checks that the harness executes. These jobs
are release preflight. They are not substitutes for the controlled corpus,
fuzzing, fault-injection, fidelity, and performance run.

\begin{table}[t]
\centering
\caption{Provenance items required for a scientifically interpretable run.}
\label{tab:reproducibility-register}
\small
\begin{tabularx}{\textwidth}{L{0.24\textwidth}Y L{0.25\textwidth}}
\toprule
Item & Required record & Integrity anchor \\
\midrule
Source and build & Commit, clean-tree state, toolchain, target, features, panic profile & Git identifier and lockfile hash \\
Policy and registries & Release policy, fidelity policy, classifier aliases, canonical registry version & SHA-256 for each frozen input \\
Corpus & Sample identifier, stratum, origin, size, expected treatment, and path & Manifest plus per-sample SHA-256 \\
External tools & FFmpeg, FFprobe, independent parsers, metric tools, and adapters & Version output and executable hash \\
Host state & Operating system, CPU, memory, storage, affinity, governor, and command line & Run manifest and environment record \\
Observations & Per-sample results, accepted outputs, raw benchmark records, coverage, fuzz logs, and faults & Schema validation and checksum inventory \\
Analysis & Deterministic scripts, thresholds, session grouping, and generated tables & Script hashes and analysis-summary schema \\
\bottomrule
\end{tabularx}
\end{table}

The controlled runner is restricted to a Linux validation server. It requires
an explicit CPU affinity, records available frequency-governor state, and
executes ten independent end-to-end benchmark sessions. Corpus
order is deterministically permuted by session. Raw JSON Lines are retained
before deterministic reduction. The run also records tool versions, host
properties, commands, logs, accepted outputs, and a SHA-256 inventory. A
development laptop may be used for compilation and narrow static inspection,
but values obtained there do not populate manuscript result macros.

The corpus evaluator validates each manifest path, size, and digest before
calling the public pipeline. It records the concrete input evidence, handler
level, verdict, alerts, canonical report, output digest, and release-gate
failures for each sample. Every clean output is processed a second time under
its canonical name and media type. Exact byte equality is required only for
deterministic structural and lossless paths. Lossy semantic paths report
same-profile canonical closure separately. They are not described as
byte-idempotent without a pre-specified normalized comparator and drift bound.

Three external adapters are mandatory for scientific release. The differential
adapter evaluates every accepted output with at least two independently
implemented parser or decoder families. The fidelity adapter applies a
separately frozen threshold policy so usability does not silently redefine the
security verdict. The fault-injection adapter covers pre-dispatch
cancellation, expired deadlines, child-launch failure, child timeout, bounded
diagnostics, output-size termination, canonical-validation failure, and
blocked publication. Each adapter exits nonzero on incomplete coverage, tool
failure, disagreement, or threshold failure and writes its report atomically.
A tooling rehearsal that disables external evidence cannot pass the scientific
release gate.

Apple release artifacts follow a separate reproducible packaging path. The
build script compiles locked static libraries for iOS device, both simulator
architectures, Mac Catalyst, and macOS. It combines architectures with
\texttt{lipo}, creates one XCFramework, validates its property list, verifies
that the strict FFI symbol exists in every archive, and compares every bundled
header with the source header.  The script sets fixed archive epochs, remaps
workspace and package-cache paths embedded by the Rust compiler, and sorts and
canonically serializes the XCFramework property list.  The CI comparison then
detects byte drift in the checked-in binary target.  These controls address the
known path, timestamp, ordering, and serialization sources of nondeterminism.
They do not imply byte-identical output across arbitrary unpinned Rust, Xcode,
SDK, or archive-tool versions, which remain part of the recorded build
environment.

Workspace publication is ordered because the facade depends on five internal
crates. The leaf crates are published before \texttt{file\_defender}, with
registry-index propagation allowed between steps. Live publication requires a
registry token, while dry run uses workspace packaging for the first release
where internal crates may not yet exist in the registry. Dependency advisories
and license/source policy are checked immediately before publication. These
controls establish release provenance and packaging consistency. They do not
establish absence of vulnerabilities in dependencies.

Numerical result macros in the master document remain \texttt{pending} until a
run manifest exists. Every frozen run is registered and retained whether it
passes or fails. A rerun receives a new identifier and links its predecessor and
reason. Only release claims require every scientific release gate to pass. Tests,
coverage-guided fuzzing, fault injection, and benchmarks are intentionally
deferred to that controlled server run. Consequently, this appendix specifies
commands, artifacts, and decision rules but reports no latency, throughput,
memory, neutralization, false-blocking, or fidelity values. This boundary
prevents local rehearsal output from being mistaken for evidence and keeps all
claims tied to a reviewable provenance record.
